% !PRINTED 56
\weberpage{56}

If now $\varkappa = 1$, then by the definition of $g$ the numbers

\[
1,\ g,\ g^2 \ldots g^{q-2}
\]

are all different from one another with respect to the modulus $q$, while $g^{q-1}$ is again congruent to 1.

But if $\varkappa > 1$, then:

\weq{4}{g^{q-1} = 1 + hq,}

and by our assumption about $g$, $h$ is not divisible by $q$. If we raise the equation (4) to the power $q$, and apply the binomial theorem on the right, then it follows that

\[
g^{q(q-1)} = 1 + hq^2 + h^2\, \frac{q^3(q-1)}{2} + \cdots,
\]

hence

\weq{5}{g^{q(q-1)} = 1 + h_1 q^2,}

wherein

\[
h_1 = h + h^2\, \frac{q(q-1)}{2} + \cdots
\]

is not divisible by $q$. (This would no longer be correct for $q = 2$, and for that reason the prime 2 requires a different treatment.)

If we raise (5) once more to the $q$th power, it follows that

\[
g^{q^2(q-1)} = 1 + h_2 q^3,
\]

and we can continue in this way, and obtain for every arbitrary positive exponent $\lambda$

\weq{6}{g^{q^{\lambda-1}(q-1)} = 1 + hq^\lambda,}

wherein $h$ is an integer not divisible by $q$. If we set for brevity:

\weq{7}{q^{\varkappa-1}(q-1) = \varphi(q^\varkappa) = c,}

then we have

\weq{8}{g^c \equiv 1\ (\mathrm{mod}\ q^\varkappa),}

and it remains to be shown that $c$ is the least positive exponent for which the congruence (8) is satisfied. Let us assume that $e$ is this least exponent, so that

\weq{9}{g^e \equiv 1\ (\mathrm{mod}\ q^\varkappa),}

then $e$ must be a divisor of $c$, because otherwise the congruence (8) would also be satisfied if $c$ were replaced by the residue of the division of $c$ by $e$, which is smaller than $e$.

On the other hand, $e$ must be divisible by $q-1$, because the congruence $g^e \equiv 1\ (\mathrm{mod}\ q)$ contained in (9) holds only for such exponents as are divisible by $q-1$.
